\documentclass[11pt,a4paper]{article}
\usepackage[a4paper,margin=23mm]{geometry}
\usepackage{lmodern,amsmath,amssymb,xcolor,booktabs}
\SetMathAlphabet{\mathbf}{normal}{OT1}{cmr}{bx}{n}
\definecolor{Definition}{HTML}{3C68A2}
\definecolor{Proof}{HTML}{9C5608}
\definecolor{Result}{HTML}{28714F}
\definecolor{Example}{HTML}{12666A}
\definecolor{Remark}{HTML}{5F6267}
\definecolor{Correction}{HTML}{A03D2C}
\usepackage[colorlinks,linkcolor=Definition,urlcolor=Definition,
pdftitle={MTL 712 - Post-Quiz 2 Notes},
pdfsubject={Finite differences for elliptic and hyperbolic PDEs and heat-equation ADI schemes}]{hyperref}

\numberwithin{equation}{section}
\setlength{\parindent}{0pt}
\setlength{\parskip}{5pt plus 1pt}
\setlength{\emergencystretch}{2em}
\setcounter{tocdepth}{2}
\raggedbottom
\pagestyle{plain}

\newsavebox{\notesbody}
\newenvironment{notesbox}[2]{%
  \par\addvspace{8pt}%
  \def\notescolor{#1}\def\notestitle{#2}%
  \setlength{\fboxrule}{0.55pt}%
  \begin{lrbox}{\notesbody}%
  \begin{minipage}{\dimexpr\linewidth-2\fboxrule-14pt\relax}%
  \setlength{\parskip}{5pt}%
  \ignorespaces
}{%
  \end{minipage}%
  \end{lrbox}%
  \noindent
  \begingroup
  \setlength{\fboxsep}{0pt}%
  \setlength{\fboxrule}{0.55pt}%
  \fcolorbox{\notescolor!68!white}{white}{%
    \begin{minipage}{\dimexpr\linewidth-2\fboxrule\relax}%
    \setlength{\fboxsep}{5pt}%
    \colorbox{\notescolor}{%
      \parbox{\dimexpr\linewidth-2\fboxsep\relax}{%
      \color{white}\small\bfseries\strut\notestitle\strut}}%
    \par\nointerlineskip
    \setlength{\fboxsep}{7pt}%
    \colorbox{\notescolor!6!white}{\usebox{\notesbody}}%
    \end{minipage}}%
  \endgroup
  \par\addvspace{8pt}%
}
\newenvironment{proofbox}[2]{%
  \begin{notesbox}{Proof}{\hypertarget{proof-#1}{\(\star\)}\enspace PROOF TO KNOW #1. #2}%
}{\end{notesbox}}
\newcommand{\prooflink}[2]{\item[\(\star\)]\hyperlink{proof-#1}{#1. #2}}
\newcommand{\step}[2]{\par\smallskip\textbf{Step #1. #2}\enspace\ignorespaces}
\newcommand{\sourcecheck}[1]{\begin{notesbox}{Correction}{About your notes}#1\end{notesbox}}
\newcommand{\datebar}[1]{{\small\itshape #1}\par}
\newcommand{\dx}{\delta_1^2}
\newcommand{\dy}{\delta_2^2}
\newcommand{\dt}{\delta_t^2}
\newcommand{\R}{\Omega}
\newcommand{\Pstep}{\tau}
\newcommand{\kron}{\otimes}
\newcommand{\bvec}{\mathbf b}
\newcommand{\uv}{\mathbf U}
\newcommand{\ones}{\mathbf 1}
\newcommand{\diag}{\operatorname{diag}}

\begin{document}

\begin{center}
\vspace*{10mm}
{\LARGE\bfseries Post-Quiz 2 Notes\par}
\smallskip
{\large MTL 712 \enspace Computational Methods for Differential Equations\par}
\smallskip
Elliptic problems, alternating-direction methods and the wave equation
\end{center}

\begin{notesbox}{Proof}{How to use these notes}
Start here after Section 20 of \texttt{till-quiz2.pdf}. These notes bring together Aditya's and Soumya's pages after the quiz on 2 November. The dated lectures span 3--16 November.

Orange boxes work through the proofs. Blue boxes give definitions and schemes. Green boxes state results, teal boxes work through examples, and red boxes explain corrections to the handwritten notes.
\end{notesbox}

\subsection*{Proofs to know}
\begin{itemize}\setlength{\itemsep}{3pt}
\prooflink{21}{Derive the five-point Laplacian and its truncation error}
\prooflink{22}{Find the eigenvalues of the Dirichlet matrix}
\prooflink{23}{Show that Peaceman--Rachford solves the steady problem}
\prooflink{24}{Prove convergence of the Dirichlet ADI iteration}
\prooflink{25}{Eliminate the stage in Douglas--Rachford}
\prooflink{26}{Find the Neumann nullspace and compatibility condition}
\prooflink{27}{Show what happens to the constant mode under ADI}
\prooflink{28}{Use the initial velocity to start the wave scheme}
\prooflink{29}{Derive the explicit wave CFL bound}
\prooflink{30}{Find the stability region of the weighted wave scheme}
\prooflink{31}{Bound the DuFort--Frankel amplification roots}
\end{itemize}

\clearpage
\begingroup
\small
\setlength{\parskip}{0pt}
\tableofcontents
\endgroup

\clearpage
\part{Elliptic equations and ADI}

\section{Elliptic equations and the Dirichlet problem}\label{sec:elliptic}
\datebar{3 November}
In two variables the principal part is \(au_{xx}+2bu_{xy}+cu_{yy}\). The equation is elliptic where \(b^2-ac<0\). The right side may depend on \(x,y,u,u_x,u_y\). We work mainly with Laplace's equation, \(\Delta u=0\), and use \(\Delta u=f\) when a source is present (Poisson's equation).

\begin{notesbox}{Definition}{The three boundary value problems}
Let \(\R=(0,1)^2\). The symbol \(\partial_n\) means differentiation along the unit normal pointing out of \(\R\).
\[
\begin{array}{@{}lll@{}}
\text{Dirichlet} & u=g & \text{prescribed values},\\
\text{Neumann} & \partial_n u=g & \text{prescribed normal derivative},\\
\text{Robin} & \alpha u+\beta\partial_n u=g & \text{a linear combination, with }\alpha,\beta>0.
\end{array}
\]
\end{notesbox}
\sourcecheck{The Robin condition in one source is missing the factor \(u\). Read its first term as \(\alpha u\).}

\subsection{The grid and the five-point stencil}
Set \(h=1/N\) and write \(U_{i,j}\) for the value at \(x_i=ih,\ y_j=jh\). There are \((N-1)^2\) interior unknowns with Dirichlet data. We insert the prescribed value \(g\) whenever the stencil reaches a boundary node.

\begin{notesbox}{Definition}{Unscaled second differences}
In space and, later, at times \(t_n=nk\), write
\[
\begin{aligned}
 \dx U_{i,j}&=U_{i+1,j}-2U_{i,j}+U_{i-1,j},\\
 \dy U_{i,j}&=U_{i,j+1}-2U_{i,j}+U_{i,j-1},\\
 \dt U^n&=U^{n+1}-2U^n+U^{n-1}.
\end{aligned}
\]
Divide spatial differences by \(h^2\) and time differences by \(k^2\) to approximate the corresponding second derivatives.
\end{notesbox}

\begin{proofbox}{21}{The five-point stencil from Taylor series}
\step{1}{Use opposite neighbours.} At \((x_i,y_j)\), expand in the \(x\)-direction.
\[
 u(x_i\pm h,y_j)=u\pm h u_x+\frac{h^2}{2}u_{xx}
 \pm\frac{h^3}{6}u_{xxx}+\frac{h^4}{24}u_{xxxx}+O(h^5).
\]
Add the two expansions and subtract \(2u\). The odd powers cancel. Keeping terms through fifth order before adding gives \(\dx u=h^2u_{xx}+\tfrac{h^4}{12}u_{xxxx}+O(h^6)\).

\step{2}{Add the two directions.} Repeat along \(y\) and divide by \(h^2\).
\[
 \Delta_h u=\frac{(\dx+\dy)u}{h^2}
 =\Delta u+\frac{h^2}{12}(u_{xxxx}+u_{yyyy})+O(h^4).
\]
The divided stencil has \(O(h^2)\) truncation error. Multiplying it by \(h^2\) gives an undivided residual of \(O(h^4)\).

\step{3}{Write the interior equation.} Replace \(\Delta u=f\) by \(\Delta_h U=f\).
\begin{equation}
 U_{i-1,j}+U_{i+1,j}+U_{i,j-1}+U_{i,j+1}-4U_{i,j}=h^2f_{i,j}.
 \label{eq:five}
\end{equation}
\end{proofbox}

\subsection{Rectangular grids and unequal steps}
\begin{notesbox}{Example}{Worked exercise. The rectangular-grid stencil}
Let \(Q=h_1/h_2\). Start from \((\dx/h_1^2+\dy/h_2^2)U=f\). Multiplying by \(h_1h_2\) gives the \(x\)-neighbours weight \(Q^{-1}\) and the \(y\)-neighbours weight \(Q\).
\[
 Q^{-1}(U_{i-1,j}+U_{i+1,j})+Q(U_{i,j-1}+U_{i,j+1})
 -2(Q+Q^{-1})U_{i,j}=h_1h_2 f_{i,j}.
\]
In the divided operator, the leading residual is
\(\tfrac1{12}(h_1^2u_{xxxx}+h_2^2u_{yyyy})\).
Multiply it by \(h_1h_2\) for the displayed stencil. For \(h_1=h_2\), the weights reduce to those in \eqref{eq:five}.
\end{notesbox}
\sourcecheck{Use fourth derivatives in the error term for the rectangular stencil. The second derivatives belong to the leading approximation to \(\Delta u\).}

\subsection{From stencil to matrix}
Multiply \eqref{eq:five} by \(-1\) to get diagonal \(4\) and neighbour coefficients \(-1\). Put the interior values into \(\uv\) with \(i\) varying fastest. Let \(L=N-1\), \(T_D=\operatorname{tridiag}(-1,2,-1)\) of size \(L\), and \(I_L\) be the identity. Then
\begin{equation}
 A\uv=\bvec,\qquad A=H+V,\qquad H=I_L\kron T_D,\qquad V=T_D\kron I_L.
 \label{eq:dir}
\end{equation}
Here \(b_{i,j}=-h^2 f_{i,j}\) plus the sum of the prescribed values at the boundary neighbours. With \(B=2I_L+T_D\), the block and stencil views are
\[
 A=\begin{pmatrix}
 B&-I_L\\
 -I_L&B&\ddots\\
 &\ddots&\ddots&-I_L\\
 &&-I_L&B
 \end{pmatrix},\quad
 B=\begin{pmatrix}
 4&-1\\
 -1&4&\ddots\\
 &\ddots&\ddots&-1\\
 &&-1&4
 \end{pmatrix},\quad
 \begin{array}{ccc}
 & -1 &\\
 -1 & 4 & -1\\
 & -1 &
 \end{array}.
\]

\begin{notesbox}{Example}{Worked example. The first interior row}
Take \(N\ge3\). At \((i,j)=(1,1)\), the left and lower neighbours lie on the boundary. Move their terms to the right to obtain
\[
 4U_{1,1}-U_{2,1}-U_{1,2}=g(0,h)+g(h,0)-h^2 f(h,h).
\]
For \(N=3\) and \(f=0\), the left side is the first row of a system with four interior unknowns.
\end{notesbox}

\subsection{Eigenvalues of the Dirichlet matrix}
\begin{proofbox}{22}{Sine modes diagonalise the five-point matrix}
\step{1}{Start with one grid line.} Take \(s_p(i)=\sin(p\pi i/N)\), \(1\le p\le N-1\). Its zero values at \(i=0,N\) fit the two ends of \(T_D\).
The sine addition formula gives
\[
 2s_p(i)-s_p(i-1)-s_p(i+1)
 =\left(2-2\cos\frac{p\pi}{N}\right)s_p(i)
 =\underbrace{4\sin^2\!\left(\frac{p\pi}{2N}\right)}_{\lambda_p}s_p(i).
\]

\step{2}{Form products.} On \(w_{pq}(i,j)=s_p(i)s_q(j)\), \(H\) acts only on the first factor and \(V\) only on the second. Hence
\[
 Hw_{pq}=\lambda_p w_{pq},\qquad Vw_{pq}=\lambda_q w_{pq},\qquad
 Aw_{pq}=(\lambda_p+\lambda_q)w_{pq}.
\]
The products form an orthogonal basis. All sums \(\lambda_p+\lambda_q\) are positive, so \(A\) is invertible. The signs here correspond to \(A\) representing \(-h^2\Delta_h\).
\end{proofbox}

\begin{notesbox}{Remark}{Why splitting by direction helps}
For a positive parameter \(\Pstep\), solving with \(I+\Pstep H\) means solving one tridiagonal system per row of the grid. With \(I+\Pstep V\), solve one per column. The two matrices commute because their actions use different coordinates. Alternating-direction implicit (ADI) methods use this to replace a coupled solve by two sweeps.
\end{notesbox}

\clearpage
\section{Peaceman--Rachford ADI}\label{sec:pr}
\datebar{6 November \enspace Lecture 38}
\subsection{Two sweeps for the heat equation}
Start from \(u_t=u_{xx}+u_{yy}\) with homogeneous Dirichlet data. Put \(r=k/h^2\). Each sweep is implicit in one direction and evaluates the other at the available stage.

\begin{notesbox}{Definition}{Peaceman--Rachford splitting}
\begin{align}
 (I-\tfrac r2\dx)U^* &= (I+\tfrac r2\dy)U^n,\nonumber\\
 (I-\tfrac r2\dy)U^{n+1} &= (I+\tfrac r2\dx)U^*.
 \label{eq:prheat}
\end{align}
First solve along \(x\) for \(U^*\), also written \(U^{n+1/2}\) in the notes. Then solve along \(y\) for \(U^{n+1}\).
\end{notesbox}
To eliminate the stage, apply \(I-\tfrac r2\dx\) to the second sweep. The directional operators commute on the homogeneous tensor grid, so we can use the first sweep on the right.
\[
 (I-\tfrac r2\dx)(I-\tfrac r2\dy)U^{n+1}
 =(I+\tfrac r2\dx)(I+\tfrac r2\dy)U^n.
\]
Expanding the products gives the Crank--Nicolson equation with the extra term
\(\tfrac{r^2}{4}\dx\dy(U^{n+1}-U^n)\) on the left. This term is \(O(k^3)\) for smooth data. After division by \(k\), the time discretisation is second order. The centred differences are second order in space.
\sourcecheck{The right side of the second sweep uses \(U^*\), the output of the first sweep.}

\clearpage
\subsection{Steady iteration with boundary data}
We can use the heat sweeps to solve a time-independent problem. Take \(A\uv=\bvec\) from \eqref{eq:dir}. Set \(\Pstep=r/2>0\) and \(\gamma=1/\Pstep\). Now \(\Pstep\) is a parameter we choose to make the iteration converge faster.

\begin{proofbox}{23}{The steady iteration solves the five-point equation}
\step{1}{Put in the boundary data.} With \(H=-\dx\) and \(V=-\dy\) acting on the interior vector, the matrix sweeps are
\begin{align}
 (\gamma I+H)\uv^* &= (\gamma I-V)\uv^n+\bvec,\nonumber\\
 (\gamma I+V)\uv^{n+1} &= (\gamma I-H)\uv^*+\bvec.
 \label{eq:prsteady}
\end{align}
Each sweep uses both directions, so each gets the full right-hand side \(\bvec\).

\step{2}{Eliminate the intermediate vector.} Multiply the second equation by \(\gamma I+H\). Use the first equation on the right. The two contributions from \(\bvec\) add to \((\gamma I-H)\bvec+(\gamma I+H)\bvec=2\gamma\bvec\).
\begin{equation}
 (\gamma I+H)(\gamma I+V)\uv^{n+1}
 =(\gamma I-H)(\gamma I-V)\uv^n+2\gamma\bvec.
 \label{eq:prelim}
\end{equation}

\step{3}{Set consecutive iterates equal.} With \(\uv^{n+1}=\uv^n=\uv\), the \(\gamma^2 I\) and \(HV\) terms cancel. What remains is \(2\gamma(H+V)\uv=2\gamma\bvec\). Conversely, if \((H+V)\uv=\bvec\), both sweeps accept \(\uv\) at all stages.
\end{proofbox}

\subsection{Error iteration and the convergence rate}
\begin{proofbox}{24}{Convergence for Dirichlet data}
\step{1}{Subtract the fixed-point equation.} Set \(e^n=\uv^n-\uv\) and subtract the equation for \(\uv\) from \eqref{eq:prelim}.
\[
 e^{n+1}=T_\Pstep e^n,\qquad
 T_\Pstep=(\gamma I+V)^{-1}(\gamma I+H)^{-1}(\gamma I-H)(\gamma I-V).
\]
In the update for \(\uv^n\), the term independent of \(\uv^n\) is
\(2\gamma(\gamma I+V)^{-1}(\gamma I+H)^{-1}\bvec\). This term drops out of the error equation.

\step{2}{Track one sine mode.} Use \(w_{pq}\) from Proof 22. Let \(a=\Pstep\lambda_p\) and \(b=\Pstep\lambda_q\). Replacing \(H,V\) by their eigenvalues gives one multiplier per mode.
\begin{equation}
 g_{pq}(\Pstep)=\frac{(1-a)(1-b)}{(1+a)(1+b)}
 =1-\frac{2(a+b)}{(1+a)(1+b)}.
 \label{eq:prg}
\end{equation}

\step{3}{Bound the multiplier.} For \(a>0\), \(|1-a|<1+a\), and the same holds for \(b\). Thus \(\rho(T_\Pstep)=\max_{p,q}|g_{pq}(\Pstep)|<1\) for fixed \(\Pstep>0\). Orthogonality of the modes gives \(\|e^n\|_2\le\rho(T_\Pstep)^n\|e^0\|_2\).
\end{proofbox}

\begin{notesbox}{Result}{Convergence rate per pair of sweeps}
The logarithmic rate is \(-\log\rho(T_\Pstep)\) when \(\rho(T_\Pstep)>0\). If \(\rho(T_\Pstep)=0\), one pair of sweeps eliminates the error.
\end{notesbox}
\begin{notesbox}{Remark}{Changing the parameter}
If we use \(\Pstep_1,\ldots,\Pstep_n\), multiply the factors \(g_{pq}(\Pstep_j)\) for each mode. On a fixed grid, choosing all parameters in a fixed interval \([\Pstep_{\min},\Pstep_{\max}]\) with \(0<\Pstep_{\min}\le\Pstep_{\max}<\infty\) gives a uniform contraction. For other choices, check that the products actually tend to zero.
\end{notesbox}
\sourcecheck{Use the spectral radius of the iteration matrix \(T_\Pstep\) in the rate formula. \(H\) describes the \(x\)-direction part of the Laplacian.}

\section{Douglas--Rachford ADI}\label{sec:dr}
Use the same two-dimensional heat equation with \(r=k/h^2\). Compare the second sweep with \eqref{eq:prheat}. Here it uses \(U^*\) and a correction evaluated at the old time level.

\begin{proofbox}{25}{From the sweeps to the steady equation}
\step{1}{Write the sweeps.}
\begin{align*}
 (I-r\dx)U^* &= (I+r\dy)U^n,\\
 (I-r\dy)U^{n+1} &= U^*-r\dy U^n.
\end{align*}

\step{2}{Eliminate \(U^*\).} Multiply the second sweep by \(I-r\dx\) and substitute the first. The two \(r\dy U^n\) terms cancel.
\begin{equation}
 (I-r\dx)(I-r\dy)U^{n+1}=(I+r^2\dx\dy)U^n.
 \label{eq:drheat}
\end{equation}

\step{3}{Set \(U^{n+1}=U^n=U\).} The mixed term cancels across the equality. This leaves \((\dx+\dy)U=0\) in the interior.
\end{proofbox}

For fixed boundary data and a fixed source, the matrix form uses \(\gamma=1/r\).
\[
\begin{aligned}
 (\gamma I+H)\uv^* &= (\gamma I-V)\uv^n+\bvec,\\
 (\gamma I+V)\uv^{n+1} &=\gamma\uv^*+V\uv^n+\bvec.
\end{aligned}
\]
Subtract the fixed point to remove \(\bvec\), then apply \eqref{eq:drheat} to one sine mode. The result is
\[
 g_{pq}=\frac{1+r^2\lambda_p\lambda_q}{(1+r\lambda_p)(1+r\lambda_q)}.
\]
It lies in \((0,1)\) on the Dirichlet grid for \(r>0\). The denominator exceeds the numerator by \(r(\lambda_p+\lambda_q)>0\).

\begin{notesbox}{Remark}{Time accuracy}
For heat flow, expanding \eqref{eq:drheat} gives backward Euler with an extra \(r^2\dx\dy(U^{n+1}-U^n)\) on the left. The scheme is first order in time. Peaceman--Rachford has the Crank--Nicolson time average and is second order in time.
\end{notesbox}

\section{Neumann problem}\label{sec:neumann}
\datebar{10 and 13 November}
Solve \(\Delta u=0\) with prescribed outward derivatives \(g_L(y),g_R(y),g_B(x),g_T(x)\). The signs of the outward derivatives are \(-u_x,u_x,-u_y,u_y\) on the left, right, bottom and top. We need \(u\) at the boundary as well as in the interior, so the unknown vector now has \((N+1)^2\) entries.

\subsection{Ghost points and the matrix}
Use the five-point equation at the boundary too. For the left side, the centred outward derivative is
\[
 g_L(y_j)=\frac{U_{-1,j}-U_{1,j}}{2h},
 \qquad U_{-1,j}=U_{1,j}+2h g_L(y_j).
\]
On the other sides the same rule gives
\[
 \begin{aligned}
 U_{N+1,j}&=U_{N-1,j}+2h g_R(y_j),\\
 U_{i,-1}&=U_{i,1}+2h g_B(x_i),\\
 U_{i,N+1}&=U_{i,N-1}+2h g_T(x_i).
 \end{aligned}
\]

\begin{notesbox}{Example}{Worked example. The lower-left corner}
Start from \(4U_{0,0}-U_{-1,0}-U_{1,0}-U_{0,-1}-U_{0,1}=0\). Substitute both ghost values. Each inward neighbour appears twice, and each side contributes one flux term.
\[
 4U_{0,0}-2U_{1,0}-2U_{0,1}=2h[g_L(0)+g_B(0)].
\]
At a non-corner side node only one direction has a doubled inward coefficient.
\end{notesbox}

Let \(T_N\) be the one-dimensional matrix of order \(N+1\).
\[
 T_N=\begin{pmatrix}
 2&-2\\
 -1&2&-1\\
 &\ddots&\ddots&\ddots\\
 &&-1&2&-1\\
 &&&-2&2
 \end{pmatrix}.
\]
With \(i\) varying fastest, set \(X=I_{N+1}\kron T_N\), \(Y=T_N\kron I_{N+1}\), and \(A=X+Y\). Then \(A\uv=2hG\). To assemble \(G\), add the outward flux from every side that touches the node. At an interior node \(G=0\).

In the block form of \(A\), all diagonal blocks are \(2I_{N+1}+T_N\). Between adjacent block rows the entries are \(-I_{N+1}\), except the inward entries of the first and last block rows, which are \(-2I_{N+1}\).

\subsection{Eigenvalues, compatibility and the constant mode}
\begin{proofbox}{26}{Cosine modes and the Neumann nullspace}
\step{1}{Fit the reflected ends.} Take \(c_p(i)=\cos(p\pi i/N)\), \(0\le p\le N\). At \(i=0\), \(2c_p(0)-2c_p(1)=2-2\cos(p\pi/N)\). The row at \(i=N\) also fits this eigenvalue. The interior rows give the same factor by the cosine addition formula.
\[
 T_Nc_p=\mu_p c_p,\qquad
 \mu_p=4\sin^2(p\pi/(2N)).
\]
The product mode \(c_p(i)c_q(j)\) has eigenvalue \(\mu_p+\mu_q\) for \(A\). The only zero eigenvalue is at \((p,q)=(0,0)\), whose vector is constant.

\step{2}{Find the left null vector.} Put \(W=\diag(\tfrac12,1,\ldots,1,\tfrac12)\) and \(W_2=W\kron W\). Multiplying the end rows by \(\tfrac12\) makes \(WT_N\) symmetric. Thus \(A\) is symmetric in the weighted inner product \(\langle u,v\rangle_{W_2}=u^{\mathsf T}W_2v\). Multiplying \(A\uv=2hG\) on the left by \(\ones^{\mathsf T}W_2\) gives the compatibility condition \(\ones^{\mathsf T}W_2G=0\).
\end{proofbox}

\begin{notesbox}{Result}{Compatibility and uniqueness}
A discrete solution exists when the weighted total flux vanishes. It is unique once we prescribe its weighted mean, since adding a constant leaves the equations intact.
\end{notesbox}

\begin{notesbox}{Remark}{How the grid weights reproduce the flux integral}
A non-corner side node has weight \(1/2\) in \(W_2\). A corner has weight \(1/4\), applied to the two fluxes there. Multiplying by \(2h\) gives weights \(h\) and \(h/2\) for each side, just as in the trapezoidal rule for \(\int_{\partial\R}g\,ds\).
\end{notesbox}

\subsection{Peaceman--Rachford for Neumann data}
Use \eqref{eq:prsteady} with the Neumann directional matrices and their right-hand side.
\[
\begin{aligned}
 (\gamma I+X)\uv^*&=(\gamma I-Y)\uv^n+2hG,\\
 (\gamma I+Y)\uv^{n+1}&=(\gamma I-X)\uv^*+2hG.
\end{aligned}
\]
Eliminating the stage as in Proof 23 gives \(\uv^{n+1}=S_\Pstep\uv^n+d_\Pstep\), where
\begin{align*}
 S_\Pstep&=(\gamma I+Y)^{-1}(\gamma I+X)^{-1}(\gamma I-X)(\gamma I-Y),\\
 d_\Pstep&=2\gamma(\gamma I+Y)^{-1}(\gamma I+X)^{-1}(2hG).
\end{align*}
The fixed-point equation is \((X+Y)\uv=2hG\).

\begin{proofbox}{27}{Why one mode survives the Neumann iteration}
\step{1}{Use the cosine products.} Substitute \(\mu_p,\mu_q\) for the two directional matrices in \(S_\Pstep\).
\[
 g_{pq}=\frac{(1-\Pstep\mu_p)(1-\Pstep\mu_q)}
              {(1+\Pstep\mu_p)(1+\Pstep\mu_q)},\qquad 0\le p,q\le N.
\]
At \(p=q=0\) the factor is \(1\). For every other mode at least one eigenvalue is positive, so \(|g_{pq}|<1\) for fixed \(\Pstep>0\).

\step{2}{Check the effect of the forcing.} Compatibility means \(G\) has zero constant component in the weighted inner product. The inverse factors preserve the cosine modes, so \(d_\Pstep\) also has zero constant component. The iteration preserves the initial weighted mean. Its error in the other modes contracts.
\end{proofbox}
\sourcecheck{There are \(N\) intervals and \(N+1\) nodes, so the sine argument in \(\mu_p\) has denominator \(2N\). The cosine products are orthogonal with the trapezoidal weights. The full iteration has spectral radius \(1\). The strict bound in the notes applies on the weighted mean-zero subspace.}

\section{Robin problem}\label{sec:robin}
\datebar{After the Neumann iteration in the source pages}
The Robin condition combines a boundary value and its normal derivative. For \(\Delta u=0\), the lecture writes the four conditions using derivatives in the positive coordinate directions.
\[
\begin{aligned}
 u_x-p_1u&=f_0(y) &&(x=0), & u_x+q_1u&=f_1(y) &&(x=1),\\
 u_y-p_2u&=g_0(x) &&(y=0), & u_y+q_2u&=g_1(x) &&(y=1).
\end{aligned}
\]
In outward form \(\partial_n u+\alpha u=\widetilde g\), the data on the left, right, bottom and top are \(-f_0,f_1,-g_0,g_1\).

\begin{notesbox}{Example}{Worked example. A Robin row on the left side}
The centred condition gives
\[
 \frac{U_{1,j}-U_{-1,j}}{2h}-p_1U_{0,j}=f_0(y_j)
 \quad\Longrightarrow\quad
 U_{-1,j}=U_{1,j}-2h[p_1U_{0,j}+f_0(y_j)].
\]
Insert this into the \(x\)-part \(2U_{0,j}-U_{1,j}-U_{-1,j}\). Move the known term to the right and then divide the row by two. The \(x\)-part is \((1+hp_1)U_{0,j}-U_{1,j}\), with right-hand contribution \(-hf_0(y_j)\).
\end{notesbox}
The other endpoint diagonals are \(1+hq_1\), \(1+hp_2\) and \(1+hq_2\). At a corner include the contribution from each incident side.

\subsection{The Robin matrix}
In one direction the matrix is
\[
 E(p,q)=\tfrac12 T_N+h\,\diag(p,0,\ldots,0,q).
\]
The interior diagonal is \(1\) and the interior off-diagonals are \(-\tfrac12\). The inward coefficient at either endpoint is \(-1\). Taking \(p=q=0\) gives the half-scaled Neumann matrix.

\clearpage
\subsection{Peaceman--Rachford for Robin data}
Set \(H_1=I_{N+1}\kron E(p_1,q_1)\) and \(V_1=E(p_2,q_2)\kron I_{N+1}\). Here \((H_1+V_1)\uv=hG\), with \(G\) the sum of outward data from the incident sides. The sweeps for that scaling are
\[
\begin{aligned}
 (\gamma I+H_1)\uv^*&=(\gamma I-V_1)\uv^n+hG,\\
 (\gamma I+V_1)\uv^{n+1}&=(\gamma I-H_1)\uv^*+hG.
\end{aligned}
\]
The same elimination as in Proof 23 gives the desired fixed-point equation.

\begin{notesbox}{Result}{Robin data remove the constant null mode}
Suppose \(p_1,q_1,p_2,q_2\ge0\) and at least one coefficient is positive. The weighted Neumann energy is a sum of squared differences between adjacent nodes. The Robin terms add nonnegative boundary squares. A vector with zero energy must be constant, and a positive boundary coefficient forces that constant to be zero. Thus \(H_1+V_1\) is positive definite in the weighted inner product.

The two directional matrices still share a weighted orthogonal basis. Each mode has at least one positive directional eigenvalue, so the bound from \eqref{eq:prg} gives convergence for fixed \(\Pstep>0\).
\end{notesbox}
\sourcecheck{Ghost elimination gives inward entry \(-1\) next to the endpoint diagonal \(1+hp\) or \(1+hq\). In the ADI equations the implicit directional matrix has a plus sign and the explicit one has a minus sign.}

\clearpage
\part{Hyperbolic equations}

\section{The wave equation in one space dimension}\label{sec:wave1}
\datebar{16 November}
Take the unit-speed wave equation on \(0<x<1\).
\[
 u_{tt}=u_{xx},\qquad u(x,0)=f(x),\qquad u_t(x,0)=g(x),
 \qquad u(0,t)=u(1,t)=0.
\]
We need initial displacement \(f\) and initial velocity \(g\), since the equation is second order in time. Put \(\nu=k/h\).

\begin{notesbox}{Definition}{Explicit three-level scheme}
Use \(\dt U^n/k^2=\dx U^n/h^2\) and solve for \(U^{n+1}\).
\begin{equation}
 U_i^{n+1}=2(1-\nu^2)U_i^n+\nu^2(U_{i-1}^n+U_{i+1}^n)-U_i^{n-1}.
 \label{eq:explicitwave}
\end{equation}
The recurrence uses two known time levels. We know \(U_i^0=f(x_i)\) and use \(g(x_i)\) to obtain \(U_i^1\).
\end{notesbox}

\subsection{First time level from the velocity}
\begin{proofbox}{28}{Starting a three-level method}
\step{1}{Centre the initial velocity.}
Write \(g_i=(U_i^1-U_i^{-1})/(2k)\) in the discrete equations. The fictitious value \(U_i^{-1}\) is \(U_i^1-2kg_i\).

\step{2}{Eliminate the fictitious value.} Set \(n=0\) in the equation \(\dt U^n=\nu^2\dx U^n\) and use \(U^0=f\).
\[
 \begin{aligned}
 U_i^1-2f_i+(U_i^1-2kg_i)&=\nu^2(f_{i-1}-2f_i+f_{i+1}),\\
 U_i^1&=f_i+kg_i+\tfrac{\nu^2}{2}(f_{i-1}-2f_i+f_{i+1}).
 \end{aligned}
\]
On the boundary set \(U^1=0\). The last term at an interior node approximates \(\tfrac{k^2}{2}u_{tt}(x_i,0)=\tfrac{k^2}{2}f''(x_i)\).
\end{proofbox}
\sourcecheck{The centred velocity uses \(2k\) in the denominator. Rearranging it gives \(U_i^{-1}=U_i^1-2kg_i\).}

\subsection{Truncation error and the CFL condition}
Taylor expansion of the smooth exact solution gives the PDE-scaled residual
\[
 \frac{\dt u}{k^2}-\frac{\dx u}{h^2}
 =\frac{k^2}{12}u_{tttt}-\frac{h^2}{12}u_{xxxx}+O(k^4+h^4).
\]
Its order is \(O(k^2+h^2)\). If you start from the undivided recurrence \eqref{eq:explicitwave}, divide the residual by \(k^2\) to get this scaling.

\begin{proofbox}{29}{Von Neumann stability of the explicit wave scheme}
\step{1}{Use a Fourier mode.} Put \(U_i^n=\zeta^n\exp(\mathrm{i}\,i\phi)\), where \(\mathrm{i}=\sqrt{-1}\). Each space shift contributes \(e^{\pm\mathrm{i}\phi}\). Thus \(\dx\) contributes \(2\cos\phi-2=-4\sin^2(\phi/2)\), and \(\dt\) contributes \(\zeta-2+\zeta^{-1}\).

\step{2}{Find the roots.} The mode equation is \(\zeta^2-2C\zeta+1=0\), with
\[
 C=1-2\nu^2\sin^2(\phi/2).
\]
The product of the roots is \(1\). For \(|C|<1\) they are conjugates on the unit circle. For \(|C|>1\) one real root has modulus greater than \(1\).

\step{3}{Check all spatial modes.} The inequality \(C\le1\) holds for every angle. The other one, \(C\ge-1\), requires \(\nu^2\sin^2(\phi/2)\le1\). The largest possible \(\sin^2(\phi/2)\) is \(1\). Therefore the mesh-independent bound is \(0<\nu\le1\). The finite-grid Dirichlet sine angles \(p\pi/N\), \(1\le p<N\), are strictly between \(0\) and \(\pi\), so the roots are distinct even at \(\nu=1\).
\end{proofbox}
\begin{notesbox}{Result}{Courant--Friedrichs--Lewy condition}
For the explicit unit-speed scheme in one dimension, take \(k/h\le1\).
\end{notesbox}

\subsection{The implicit theta family}
In the spatial difference, replace \(U^n\) by
\(\theta U^{n+1}+(1-2\theta)U^n+\theta U^{n-1}=U^n+\theta\dt U^n\). Move the new term to the left.
\begin{equation}
 (I-\theta\nu^2\dx)\dt U^n=\nu^2\dx U^n,\qquad \theta\ge0.
 \label{eq:weightedwave}
\end{equation}
The case \(\theta=0\) gives the explicit scheme. The lecture calls \(\theta=\tfrac14\) the von Neumann scheme. For \(\theta>0\), solve a tridiagonal system for \(U^{n+1}\). The symmetric time average changes the scaled spatial difference by \(O(k^2)\), so the residual remains \(O(k^2+h^2)\).

\begin{proofbox}{30}{Stability of the weighted wave scheme}
\step{1}{Substitute the same Fourier mode.} Write \(s=\sin^2(\phi/2)\). The amplification equation becomes
\[
 (1+4\theta\nu^2s)(\zeta-2+\zeta^{-1})=-4\nu^2s,\qquad
 \frac{\zeta+\zeta^{-1}}{2}=1-\frac{2\nu^2s}{1+4\theta\nu^2s}.
\]

\step{2}{Apply the root bound from Proof 29.} The right side is at most \(1\). Its lower bound of \(-1\) is equivalent to
\[
 (1-4\theta)\nu^2 s\le1.
\]
For \(\theta\ge\tfrac14\) this holds for all \(\nu\). For \(0\le\theta<\tfrac14\), use \(s\le1\) and require \(\nu^2\le1/(1-4\theta)\).
\end{proofbox}

\clearpage
\section{The wave equation in two space dimensions}\label{sec:wave2}
Take \(u_{tt}=u_{xx}+u_{yy}\) on the unit square with \(u(x,y,0)=f(x,y)\), \(u_t(x,y,0)=g(x,y)\), and \(u=0\) along the boundary. Set \(D=\dx+\dy\).

\begin{notesbox}{Definition}{The 2D schemes}
The explicit stencil is \(\dt U^n=\nu^2 D U^n\). Solving for the top level gives
\[
\begin{aligned}
 U_{i,j}^{n+1}&=2(1-2\nu^2)U_{i,j}^n-U_{i,j}^{n-1}\\
 &\quad+\nu^2(U_{i-1,j}^n+U_{i+1,j}^n+U_{i,j-1}^n+U_{i,j+1}^n).
\end{aligned}
\]
The weighted method is \((I-\theta\nu^2D)\dt U^n=\nu^2 D U^n\). Both have \(O(k^2+h^2)\) scaled truncation error under the same smoothness assumptions as in one dimension.
\end{notesbox}

\textbf{Startup.} Use \(U^{-1}=U^1-2kg\) as before. For the weighted scheme this gives
\[
 (I-\theta\nu^2D)(U^1-f-kg)=\tfrac{\nu^2}{2}Df.
\]
With \(\theta=0\), read this as \(U^1=f+kg+\tfrac{\nu^2}{2}Df\). The same formula with \(D=\dx\) starts the one-dimensional method.

\textbf{Stability.} Each space direction contributes a term to the Fourier calculation. Replace \(s\) in Proof 30 by \(\sin^2(\phi/2)+\sin^2(\psi/2)\), whose maximum is \(2\). This is the only change in the modal bound.

\begin{notesbox}{Result}{Wave stability on uniform grids}
\begin{center}
\begin{tabular}{@{}lll@{}}
\toprule
Scheme & One space dimension & Two space dimensions\\
\midrule
Explicit, \(\theta=0\) & \(\nu\le1\) & \(\nu\le1/\sqrt{2}\)\\[2pt]
\(0<\theta<\tfrac14\) & \(\nu^2\le1/(1-4\theta)\) & \(\nu^2\le1/[2(1-4\theta)]\)\\[2pt]
\(\theta\ge\tfrac14\) & no restriction on \(\nu\) & no restriction on \(\nu\)\\
\bottomrule
\end{tabular}
\end{center}
Here \(\nu=k/h\). These are the mesh-independent bounds from the mode calculation.
\end{notesbox}

\sourcecheck{The square problem needs boundary values along all four sides for the sine-mode calculation. One source lists only corner values. The notes mention a possible wave-equation ADI scheme and then switch to the heat-equation exercise.}

\clearpage
\part{Heat-equation exercise}

\section{A single sine mode under four schemes}\label{sec:exercise}
\datebar{Exercise following the wave lecture}
\begin{notesbox}{Example}{Set up the exercise}
Solve \(u_t=u_{xx}\) on \(0<x<1\), with \(u(x,0)=\sin(\pi x)\) and \(u(0,t)=u(1,t)=0\). The exact solution is \(e^{-\pi^2t}\sin(\pi x)\).

On the grid write \(s_i=\sin(\pi i/N)\) and \(U_i^n=a_n s_i\). Set \(r=k/h^2\), \(\phi=\pi/N\), and \(\sigma=\sin^2(\phi/2)\). Since \(\dx s_i=-4\sigma s_i\), each new time level stays in the same spatial mode. The problem is to find the scalar \(a_n\).
\end{notesbox}

\subsection{Schmidt, Crank--Nicolson and Douglas}
Each of these schemes relates two time levels.
\begin{align*}
 U^{n+1} &=(I+r\dx)U^n &&\text{(Schmidt / explicit)},\\
 (I-\tfrac r2\dx)U^{n+1}&=(I+\tfrac r2\dx)U^n &&\text{(Crank--Nicolson)},\\
 \bigl[I-\tfrac12(r-\tfrac16)\dx\bigr]U^{n+1}
 &=\bigl[I+\tfrac12(r+\tfrac16)\dx\bigr]U^n &&\text{(Douglas compact)}.
\end{align*}

\begin{notesbox}{Example}{Worked solution. Read off the growth factor}
For Crank--Nicolson, substitute \(a_n s_i\) and cancel the common spatial mode. This gives \((1+2r\sigma)a_{n+1}=(1-2r\sigma)a_n\). Divide by \(1+2r\sigma\). The other rows follow by the same substitution. Since \(a_0=1\), in each case \(a_n=G^n\).
\begin{center}
\renewcommand{\arraystretch}{1.5}
\begin{tabular}{@{}ll@{}}
\toprule
Method & \(G=a_{n+1}/a_n\) when \(a_n\ne0\)\\
\midrule
Schmidt & \(1-4r\sigma\)\\
Crank--Nicolson & \(\dfrac{1-2r\sigma}{1+2r\sigma}\)\\
Douglas compact & \(\dfrac{1-2(r+\tfrac16)\sigma}{1+2(r-\tfrac16)\sigma}\)\\
\bottomrule
\end{tabular}
\end{center}
\end{notesbox}

\textbf{Stability across modes.} Replace \(\sigma\) by \(\sin^2(p\pi/(2N))\). The usual full-mode bound for Schmidt is \(r\le\tfrac12\). Crank--Nicolson and Douglas have modulus at most \(1\) for every \(r>0\). In particular, write the Douglas quotient as \((a-b)/(a+b)\), where \(a=1-\sigma/3>0\) and \(b=2r\sigma\ge0\).

\clearpage
\subsection{DuFort--Frankel}
Centre the time derivative and replace the centre value in the spatial second difference by \((U_i^{n+1}+U_i^{n-1})/2\). Gather the terms at time \(n+1\).
\[
 (1+2r)U_i^{n+1}=2r(U_{i-1}^n+U_{i+1}^n)+(1-2r)U_i^{n-1}.
\]

\begin{notesbox}{Example}{Worked solution. A recurrence for the same sine mode}
Substitute \(U_i^n=a_n s_i\). The two neighbours sum to \(2a_n\cos(\phi)s_i\). Hence
\[
 (1+2r)a_{n+1}=4r\cos(\phi)a_n+(1-2r)a_{n-1}.
\]
Start with \(a_0=1\) and take one Crank--Nicolson step \(a_1=(1-2r\sigma)/(1+2r\sigma)\), for example. Use the recurrence to advance the remaining levels.
\end{notesbox}

\subsection{DuFort--Frankel stability and consistency}
The undated fragment above the 3 November heading contains the end of this stability calculation.

\begin{proofbox}{31}{The two DuFort--Frankel amplification roots}
\step{1}{Solve the modal recurrence.} Use an arbitrary Fourier angle \(\phi\) and write \(a_n=\zeta^n\). The scalar quadratic has roots
\[
 \zeta_\pm=\frac{2r\cos\phi\ \pm\sqrt{1-4r^2\sin^2\phi}}{1+2r}.
\]

\step{2}{Check the real-root case.} When \(1-4r^2\sin^2\phi\ge0\), the square root is at most \(1\). The triangle inequality gives
\[
 |\zeta_\pm|\le\frac{2r+1}{1+2r}=1.
\]
If the discriminant is zero, the repeated root has modulus less than \(1\).

\step{3}{Check the complex-root case.} The roots are conjugates. Their squared modulus is
\[
 |\zeta_\pm|^2
 =\frac{4r^2\cos^2\phi+4r^2\sin^2\phi-1}{(1+2r)^2}
 =\frac{2r-1}{2r+1}<1.
\]
This proves the unconditional von Neumann root bound.
\end{proofbox}

\begin{notesbox}{Result}{Consistency requires \(k/h\to0\)}
The replacement of \(2u(t)\) by \(u(t+k)+u(t-k)\) adds a term \((k/h)^2u_{tt}\) to the PDE-scaled residual. The heat equation is recovered for smooth general solutions as \(h,k\to0\) with \(k/h\to0\). The root bound holds for all \(r>0\).
\end{notesbox}

\clearpage
\subsection{The two-dimensional heat formulas}
The last lecture pages give the Peaceman--Rachford product from Section~\ref{sec:pr} and a second product using the one-dimensional Douglas compact factors. For \(j=1,2\), write
\[
 L_j=I-\tfrac12(r-\tfrac16)\delta_j^2,\qquad
 R_j=I+\tfrac12(r+\tfrac16)\delta_j^2.
\]
The compact equation is \(L_1L_2U^{n+1}=R_1R_2U^n\). To implement it, solve
\[
 L_1U^*=R_2U^n,\qquad L_2U^{n+1}=R_1U^*.
\]
Multiplying the second sweep by \(L_1\) and commuting it with \(R_1\) recovers the product equation.

\begin{notesbox}{Remark}{Where the compact correction comes from}
In one direction the scheme is Crank--Nicolson with spatial operator \(h^{-2}(I+\delta_j^2/12)^{-1}\delta_j^2\). The correction \(I+\delta_j^2/12\) cancels the leading \(h^2\partial_j^4u/12\) term in the Taylor expansion of the ordinary centred second difference.
\end{notesbox}

\sourcecheck{On the final pages, the two sweeps labelled Peaceman--Rachford have the Douglas--Rachford form in Section~\ref{sec:dr}. The last page labels the compact product Douglas--Rachford and mixes the stage indices. The compact sweeps above reproduce its product equation.}

\section*{Sources and lecture coverage}
\phantomsection
\addcontentsline{toc}{section}{Sources and lecture coverage}

\begingroup\small
\begin{tabular}{@{}ll@{}}
\toprule
Date & Material in the supplied notes\\
\midrule
3 Nov. & Elliptic equations and the Dirichlet five-point stencil\\
6 Nov. & Peaceman--Rachford ADI (lecture 38)\\
10 and 13 Nov. & Neumann problem and its iteration, followed by Robin data\\
16 Nov. & Wave equation, followed by a heat-equation exercise\\
\bottomrule
\end{tabular}

\begin{itemize}\setlength{\itemsep}{2pt}\setlength{\parsep}{0pt}
\item \texttt{till-quiz2.pdf}, pp.\ 36--38, ends with the explicit two-dimensional heat scheme and its stability.
\item Soumya's handwritten PDF, \emph{MTL712 (Computational Methods For Differential Equations) 5}, has the elliptic material on pp.\ 1--7 and the wave and heat material on pp.\ 8--10. The DuFort--Frankel fragment at the top of p.\ 1 is undated. The accompanying notebook photos from 19:13:52 to 19:13:54 include the 6 November ADI lecture.
\item Aditya's notebook photos supply the Douglas--Rachford factorisation and fill in matrix and convergence steps. They also cover the wave schemes and the four-method heat exercise. Undated pages follow the topic order. The 19 November filenames record WhatsApp timestamps.
\item Soumya's \emph{MTL712 Elliptic PDE Combined Notes} provided a second reading of the elliptic pages.
\end{itemize}

\textbf{ADI references.}
Peaceman and Rachford, \emph{The Numerical Solution of Parabolic and Elliptic Differential Equations} (1955). \href{https://doi.org/10.1137/0103003}{Paper}

Douglas and Rachford, \emph{On the Numerical Solution of Heat Conduction Problems in Two and Three Space Variables} (1956). \href{https://doi.org/10.1090/S0002-9947-1956-0084194-4}{Paper}
\endgroup

\end{document}
